\ifJA

\section*{はじめに}
\addcontentsline{toc}{section}{はじめに}

『伊勢物語』には、長いあいだ多くの注釈と現代語訳が積み重ねられてきました。
それらは、分かりにくいことばを説明し、人物や出来事の関係を明らかにし、
物語を理解するための多くの知識を与えてくれます。

しかし、説明を加えるほど、本文そのものが許していた読み方を
一つに決めてしまうこともあります。

この本では、できるだけ逆のことをしてみます。

本文が決めていないことは、急いで決めない。
分からないことは、分からないまま残す。
複数に読めるところは、なるべく複数に読めるままにする。

それは、何をどう読んでもよいという意味ではありません。
本文に書かれている人、物、場所、時間、行動をできるだけそのまま残し、
そこから読める範囲を狭めない、ということです。

題名を『読み手勝手伊勢物語』としました。

訳者が勝手に決めないから、読み手が勝手に読める。

そのための本です。

\else

\section*{Preface}
\addcontentsline{toc}{section}{Preface}

Over the centuries, \textit{The Tale of Ise} has accumulated a rich tradition
of commentary and translation. These works explain difficult words,
clarify relations among people and events, and provide knowledge that helps
us understand the text.

Yet the more we explain, the more easily we may settle something
that the text itself has left unsettled.

This book tries to move in the opposite direction.

What the text does not determine, we do not rush to determine.
What remains unclear may remain unclear.
Where more than one reading is possible, we try to preserve that possibility.

This does not mean that anything goes.
People, objects, places, times, and actions given in the text are kept in place
as far as possible, without unnecessarily narrowing what can be read from them.

The Japanese title of this book is
\textit{Yomite Katte Ise Monogatari}.
Its English title is \textit{The Tale of Ise --- BYO Reading}.

The translator decides less,
so that the reader can read more.

That is what this book is for.

\fi
